\documentclass[11pt]{article}
\usepackage[margin=1in]{geometry}
\usepackage{amsmath,amssymb,bm,booktabs,graphicx,hyperref}
\hypersetup{colorlinks=true,linkcolor=blue,urlcolor=blue}

\title{CS 577 Deep Learning\\Homework 2: Inside an MLP}
\author{Your Name (Hawk username: \texttt{yourname})}
\date{Fall 2026}

\begin{document}
\maketitle

\section*{Part 0: Conventions and Reproducible Setup}

\subsection*{0.1 Environment and public tests}
% Report versions and summarize the public-test result.

\subsection*{0.2 Shape and reduction contract}
% Explain (B,1) versus (B,), the possible (B,B) broadcasting error, and 1/B.

\section*{Part 1: Forward and Backward Mathematics}

\subsection*{1.1 Forward shape trace}
\begin{center}
  \begin{tabular}{lll}
    \toprule
    Quantity & Expression / role & Shape \\
    \midrule
    $X$ & input batch & TODO \\
    $W_1,b_1$ & first affine parameters & TODO \\
    $Z_1$ & $XW_1+b_1$ & TODO \\
    $A_1$ & $\operatorname{ReLU}(Z_1)$ & TODO \\
    $W_2,b_2$ & output affine parameters & TODO \\
    $S$ & logits & TODO \\
    $y$ & binary targets & TODO \\
    $J$ & mean batch loss & TODO \\
    \bottomrule
  \end{tabular}
\end{center}

% Explain both bias broadcasts and distinguish logits from probabilities.

\subsection*{1.2 BCE-with-logits derivative}
% Derive dJ/dS = (sigmoid(S)-y)/B with intermediate reasoning.

\subsection*{1.3 Complete backward chain}
% Derive dW2, db2, dA1, dZ1, dW1, db1, and optionally dX.
% Include shapes and explain batch-axis reductions for bias gradients.

\section*{Part 2: NumPy Forward and Backward Implementation}

\subsection*{2.1 Stable primitives}
% Explain the stability choices and report one extreme-logit check.

\subsection*{2.2 Fixed-case forward pass}
% Report Z1/A1/logit shapes and logits from the supplied check case.

\subsection*{2.3 Fixed-case manual backward pass}
\begin{center}
  \begin{tabular}{lrr}
    \toprule
    Parameter & Gradient shape & Gradient norm \\
    \midrule
    $W_1$ & TODO & TODO \\
    $b_1$ & TODO & TODO \\
    $W_2$ & TODO & TODO \\
    $b_2$ & TODO & TODO \\
    \bottomrule
  \end{tabular}
\end{center}

% Report fixed-case loss and briefly interpret the gradient norms.

\section*{Part 3: Two Independent Gradient Checks}

\subsection*{3.1 Centered finite differences}
\begin{center}
  \begin{tabular}{lrr}
    \toprule
    Parameter & Maximum absolute error & Relative error \\
    \midrule
    $W_1$ & TODO & TODO \\
    $b_1$ & TODO & TODO \\
    $W_2$ & TODO & TODO \\
    $b_2$ & TODO & TODO \\
    \bottomrule
  \end{tabular}
\end{center}

\subsection*{3.2 PyTorch autograd comparison}
% Provide a similar table for autograd and explain what passing establishes.

\section*{Part 4: Nonlinear Experiment}

\subsection*{4.1 Linear baseline}
% Report train/test loss and accuracy; explain the straight-boundary limitation.

\subsection*{4.2 Manual MLP}
% Report configuration and train/test loss and accuracy.

% Insert the loss curve and decision-boundary figure exported by the notebook.
% Example: \includegraphics[width=.8\linewidth]{hw2_experiment.png}

\subsection*{4.3 Interpretation}
% Write the requested 5--8 sentence evidence-based interpretation.

\section*{Part 5: Reflection}

\paragraph{Most useful cached value.} TODO

\paragraph{First gradient-check investigations.} TODO

\paragraph{Why close is not bit-identical.} TODO

\section*{Assistance and Generative-AI Disclosure}

% Name each tool/person, purpose, and what you verified or changed. If none, state that.
TODO

\end{document}
